\documentclass[dvipdfmx,a4paper,oneside]{jsbook}
\usepackage[left=20truemm,right=20truemm,top=25truemm,bottom=25truemm]{geometry}
\usepackage[utf8]{inputenc}
\usepackage{amsmath}
\usepackage{amsfonts}
\usepackage{amssymb}
\usepackage[dvipdfmx,pageanchor=false]{hyperref}
\usepackage{pxjahyper}
\usepackage{setspace}
\usepackage{ascmac}

\title{光はどこにいるか？ \\二つの時空と並進で紐解く相対論と量子力学の謎\\
ニュートン‐ウィグナー問題とは？（PGA版改訂）}
\author{二重時空理論プロジェクト}
\date{}

\begin{document}

\frontmatter
\maketitle
\tableofcontents

\mainmatter
\renewcommand{\sectionmark}[1]{\markboth{#1}{}}

\section{ニュートン‐ウィグナー問題とは何か？}

みなさん、こんにちは！ この読み物では、相対論と量子力学の不思議な世界を探検します。まずは、ちょっとした謎から始めましょう。

みなさんは、光（光子）は「どこにいる」と考えますか？ 普通の粒子、たとえば電子なら「ここにいる」と簡単に言えます。でも、光子のような質量がゼロの粒子になると、話が少しややこしくなります。

この問題は、70年以上にわたり物理学者を悩ませ続け、量子場論の発展やリー・シュリーダーの定理などの議論でも繰り返し取り上げられてきました。相対論と量子論の統合が、いかに深い難問を抱えているかを象徴するものです。

特殊相対性理論では、光はいつも光速で進みます。このとき、光の進行方向に沿って見ると、空間の長さがどんどん縮んで、ついには無限に「つぶれて」しまいます。イメージとしては、速すぎる電車に乗っていると、電車の長さがゼロに見えてしまうような感じです。

このせいで、光子の「位置」を正確に決めることがとても難しくなります。数学的に言うと、進行方向の位置が「つぶれて」しまい、きれいで一意な位置演算子（位置を表す量）が作れなくなってしまうのです。

\begin{itembox}[l]{ニュートン‐ウィグナー問題の起源}

この問題は、1940年代後半にまで遡ります。量子力学と特殊相対性理論を統合しようとした時代、著名な物理学者ユージン・ウィグナー（Eugene Wigner、1963年にノーベル物理学賞を受賞した群論の大家）とテオドール・ニュートン（T. D. Newton）が、1949年に重要な論文を発表しました（Localized States for Elementary Systems, Reviews of Modern Physics, 21, 400）。この論文で、彼らは相対論的量子力学における「粒子の位置演算子」の定義を厳密に検討し、質量のある粒子では位置をうまく定義できる一方、質量零の粒子（特に光子）では共変的で一意な位置演算子が存在しないことを指摘しました。これが\textbf{ニュートン‐ウィグナー問題}（Newton-Wigner localization problem）の起源です。

\end{itembox}

簡単にまとめると：
\begin{itemize}
  \item 質量のある粒子（電子など）では、位置演算子をうまく定義できます。
  \item しかし、光子では進行方向の情報が失われ、位置が光の進む面全体に広がったようにしか言えなくなります。
  \item 違う速さで動く観測者同士で、位置の値が一致しにくく、相対性理論の美しさが崩れてしまいます。
\end{itemize}

この問題は、量子力学と相対性理論を一緒に考えるときの大きな壁でした。長年、多くの物理学者を悩ませてきた謎なのです。

それでは、この謎をどのように解くのでしょうか？ 次で、二重時空理論（特にPGA版）の出番です！

\section{二重時空理論（PGA版）がどのように解決するか？}

みなさん、ニュートン‐ウィグナー問題の謎を感じていただけましたか？ ここで、二重時空理論のPGA版が登場します。この理論は、まったく新しい視点でこの問題をすっきり解決します。

二重時空理論の大きなアイデアは、「時空は一つだけじゃない！」というものです。すべての粒子は、内在的に二つの時空を持っています：
\begin{itemize}
  \item \textbf{通常時空}：私たちが普段見ている普通の時空です。ここでは、光子の進行方向が「つぶれて」しまいます。
  \item \textbf{双対時空}：各粒子が持つ、もう一つの特別な時空です。この双対時空はコンパクト化（有限でぐるぐる巻きついた）されていて、時間の矢印が逆向きです。
\end{itemize}

PGA版（射影幾何代数 \(G(3,1,1)\) への拡張）では、さらに重要な自由度が加わります。  
それは\textbf{並進}（空間や時間の実際の移動を司る自由度）です。通常のCl(3,1)に、ヌルベクトル \(e_4\)（平方がゼロの特別な方向）を加えることで、4つのヌル生成子 \(N_\mu = e_4 \wedge e_\mu\) が生まれます。これらが「実際の位置のずれ」を担います。

光子の場合、質量がゼロなので、通常時空と双対時空のローターが\textbf{完全反同期}というきれいな状態になります。この同期のおかげで：
\begin{itemize}
  \item 通常時空で「つぶれた」進行方向の情報が、双対時空のコンパクトな位相（角度のようなもの）として、ぴったりと保存されます。
  \item さらにPGA版では、ねじれ不整合がゼロに固定されると、並進パラメータ \(\lambda^\mu\) が光円錐上に強制されます（\(\eta_{\mu\nu}\lambda^\mu\lambda^\nu = 0\)）。つまり、光子の移動は純粋な「光速並進」になります。
\end{itemize}

イメージとしては、二つの時空と並進が「協力」している感じです。  
通常時空で失われた情報を双対時空が受け止め、実際の移動は並進セクターが光的に実行する。  
この三重の役割分担により、光子の位置は「通常時空の横方向 ＋ 双対時空の進行方向位相 ＋ 光的並進」という形で、点のように一意に定義できます。

位置演算子もすべての観測者で一致する美しい形になり、ニュートン‐ウィグナー問題が解決されるのです！

\section{双対時空における「潰れ」が生じない理由と、並進の役割}

みなさん、ここでは前節で出てきた「潰れ」（長さの収縮）が、双対時空ではなぜ起こらないのかを、できるだけ分かりやすくお話しします。PGA版で加わった並進の役割も一緒に見ていきましょう。

まず、通常時空での「潰れ」をもう一度おさらいしましょう。

通常時空は、私たちが普段考えている普通の時空です。光子が光速で進むとき、進行方向（たとえばz方向）にローレンツブーストという変換がかかります。これは、数学的に「双曲的な動き」（\(\hat{p}^2 = +1\)）と呼ばれ、非コンパクト（無限に伸びる）性質を持っています。

イメージとしては：
\begin{itemize}
  \item 長いゴムひもをどんどん引っ張るような感じです。引っ張れば引っ張るほど、長さが無限に伸びたり縮んだりして、進行方向の位置が「つぶれて」しまいます。
  \item 結果、光子の進行方向の自由度が失われ、位置を正確に特定できなくなります。
\end{itemize}

それでは、双対時空ではどうでしょうか？

双対時空は、各粒子が内在的に持つもう一つの時空で、コンパクト化（有限で巻きついた）された構造をしています。光子の場合、通常時空と双対時空のローターは\textbf{完全反同期}という特別な状態になります。これは、
\[
R_{\text{usual}} = \exp\left( \frac{\phi}{2} \hat{p} \right), \quad
R_{\text{dual}} = \exp\left( \frac{\theta}{2} \hat{q} \right)
\]
のように表され、\(\phi = \theta\)（角度が同じ）で、方向が逆向き（\(\hat{p} i = -\hat{q}\)）になる関係です。

ここが大事なポイントです。双対時空の動きは：
\begin{itemize}
  \item \textbf{回転型}（\(\hat{q}^2 = -1\)）で、楕円的（コンパクト）です。
  \item 位相角（角度）\(\theta\) は\(4\pi\)周期で繰り返す、輪っかのような有限の空間です。
  \item さらに、双対時空の時間基底は逆向き（時間矢印が反転）なので、通常時空の「引っ張るようなブースト」が、双対時空では「ぐるぐる回る回転」に変わります。
\end{itemize}

イメージとしては：
\begin{itemize}
  \item 輪ゴムを回すような感じです。いくら速く回しても、輪ゴムの長さは変わらず、ただ角度（位相）がどんどん蓄積されるだけです。
  \item 通常時空で「無限に縮んだ」進行方向の情報は、双対時空でこの有限の位相角としてぴったり保存されます。
\end{itemize}

PGA版で加わった\textbf{並進}の役割は、さらに重要です。  
並進はヌル生成子 \(N_\mu\) によって担われ、平方がゼロという性質を持っています。光子ではねじれ不整合がゼロに固定されるため、並進パラメータは光円錐上に強制されます。  
つまり「実際の移動」は、無限に縮む心配のない、純粋な光速シフトとして実行されるのです。

これにより、光子の位置は：
\begin{itemize}
  \item 通常時空の横方向の情報
  \item 双対時空のコンパクトな進行方向位相
  \item 光的並進による実際の変位
\end{itemize}
の三つの要素で、点のように正確に表現できるようになります。

\section{ヒルベルト空間の構成}

みなさん、ここでは量子力学の大事な舞台である「ヒルベルト空間」を、二重時空理論ではどのように作るのかをお話しします。少し抽象的になりますが、イメージをたくさん使って丁寧に説明しますので、一緒に進めていきましょう！

まず、ヒルベルト空間とは何でしょうか？

量子力学では、粒子の状態（たとえば「ここにいて、この速さで動いている」といった情報）をすべて表す「場所」が必要です。それがヒルベルト空間です。普通の量子力学では、位置や運動量の波動関数で表現しますが、二重時空理論では、少し違った美しい方法で作ります。

二重時空理論のヒルベルト空間は、主に\textbf{双対時空のコンパクトな位相空間}を基盤にしています。

双対時空のローター（回転を表す量）は、全体として\(S^3\)（3次元球の表面のようなコンパクトな空間）をぴったりとパラメトライズします。光速が一定であるため、双対時空の「半径」は厳密に固定され、方向と角度が一体となった\(S^3\)だけで完全に記述できるのです。

イメージとしては：
\begin{itemize}
  \item \(S^3\)：風船のような球の表面を想像してください。この表面全体が、双対時空のローターの方向と角度を同時に表しています。
  \item 風船の表面をぐるぐる動く点が、光子や粒子の内部状態に対応します。
\end{itemize}

数学的には、ヒルベルト空間を
\[
\mathcal{H} = L^2(S^3)
\]
のように表します。これは「\(S^3\)上の、平方可積分な波動関数全体の集まり」という意味です。

PGA版では、この構造に並進の自由度が自然に加わります。並進はヌル生成子で記述されるため、光的な場合は「どの光円錐方向に進むか」という情報が、双対時空の位相と結びついて表現されます。全体として、光子の全情報（位置、偏光、波動性、進行方向）がきれいに収まります。

このヒルベルト空間の素晴らしいところは：
\begin{itemize}
  \item 双対時空が完全にコンパクト（\(S^3\)）なので、自然に量子的な離散性（小さなスケールでの切れ味の良さ）が現れます。
  \item 通常時空のローターと並進がこの空間に作用すると、ローレンツ変換と並進が正しく保たれ、すべての観測者で状態が一貫します。
\end{itemize}

イメージとしては、「風船の表面の上で、波がゆらゆら揺れている」様子を思い浮かべてください。この波が光子や粒子の量子状態そのものなのです。

\section{位置演算子の定義とスピノル表現との統合}

みなさん、ここではヒルベルト空間に「位置演算子」を定義し、それが粒子のスピン（スピノル表現）とどのように美しくつながるのかをお話しします。少し専門的な言葉が出てきますが、イメージを大切にしながら丁寧に進めますね！

まず、位置演算子とは何でしょうか？

量子力学では、粒子の「位置」を測るための特別な量を位置演算子と呼びます。普通の量子力学では、位置演算子は波動関数に「その場所の値を掛ける」ような働きをします。これにより、「粒子がここにいる確率」を計算できます。

二重時空理論（PGA版）では、この位置演算子を\textbf{双対時空のコンパクトな空間基底}と\textbf{並進セクター}を組み合わせて定義します。

双対時空の空間部分は、
\begin{itemize}
  \item \{jI, jJ, jK\}
\end{itemize}
という基底で表され、これらはコンパクト（有限で巻きついた）な性質を持っています。  
さらにPGA版では、ヌル生成子 \(N_\mu\) が実際の変位を担います。

イメージとしては：
\begin{itemize}
  \item 双対時空の空間を、小さな輪っかや球の上で動く「座標」のように考えます。
  \item 位置演算子は、この輪っか上で「今、どの角度にいるか」を掛ける演算子に、並進による実際のシフトを加えたものです。
\end{itemize}

数学的には、位置演算子を
\[
\hat{X}' = x' jI + y' jJ + z' jK + \text{（並進によるシフト項）}
\]
のように表し、状態（波動関数）にこの値を掛けます。特に、進行方向の位置は双対時空の位相角 \(\theta\) と光的並進パラメータに依存するので、通常時空で失われた情報がここでぴったり回復されます。

この位置演算子の良いところは：
\begin{itemize}
  \item 運動量演算子（位相を微分する演算子）と正しい交換関係（\([\hat{X}, \hat{P}] = i\hbar\)）を持っています。
  \item ローレンツ変換と並進（ポアンカレ変換）に対しても、すべての観測者で一貫して振る舞います。
  \item 光子の場合、並進が光円錐に強制されるため、位置の「つぶれる」自由度が最初から存在しません。
\end{itemize}

これで、光子の位置が点のように扱え、ニュートン‐ウィグナー問題が解決されます！

次に、スピノル表現との統合についてです。

スピノルとは、粒子のスピン（回転の性質）を表す特別な量子状態です。電子などの粒子は1/2スピンを持ち、ディラック方程式で記述されます。二重時空理論では、このスピノルが双対時空のローターと自然につながります。

双対時空のローター \(R_{\text{dual}}\) は、数学的にCl(3,1)という代数のスピノル表現を生成します。具体的には：
\begin{itemize}
  \item 偽スカラー（特別な要素）\(i\) が、左巻きと右巻きのスピノル（キラル投影）を生み出します。
  \item ヒルベルト空間の特定の部分（スピン1/2の状態）が、4成分のディラックスピノルにぴったり対応します。
\end{itemize}

PGA版では、モーター \(M = R\, T\)（ローターと並進の積）が、位置とスピンを統一的に変換する役割を果たします。サンドウィッチ積 \(M X \widetilde{M}\) によって、位置と向きが同時に正しく変換されるのです。

この統合により、ディラック方程式の質量項は双対ローターの「剛性」として現れ、ねじれ不整合（通常時空と双対時空のずれ）が粒子の運動や相互作用を説明します。光子では剛性がゼロなので、純粋な光的並進だけが残ります。

\section{まとめ}

みなさん、ここまで一緒に二重時空理論（PGA版）の量子的な部分を学んできました。お疲れ様でした！ 最後に、これまでの内容を振り返りながら、全体像をおさらいしましょう。

二重時空理論は、従来の相対論と量子力学の難しい問題（特にニュートン‐ウィグナー問題）を、まったく新しい視点で解決します。その鍵は：
\begin{itemize}
  \item すべての粒子が「通常時空」と「双対時空」という二つの時空を内在的に持っていることです。
  \item 双対時空はコンパクト（有限で巻きついた）で、時間の矢印が逆向きという特別な性質があります。
  \item PGA版では、さらに並進（ヌル生成子）が加わり、実際の空間移動を内在的に記述できるようになります。
\end{itemize}

光子のような質量零粒子では：
\begin{itemize}
  \item 通常時空で進行方向が「つぶれて」位置が決めにくくなる問題がありましたが、
  \item 双対時空のコンパクトな位相（角度）がその失われた情報をぴったり保存してくれます。
  \item さらに、ねじれ不整合がゼロに固定されることで、並進が光円錐上に強制され、純粋な光速移動として実行されます。
  \item その結果、光子を点粒子として美しく扱え、位置演算子がすべての観測者で一貫するようになります。
\end{itemize}

さらに：
\begin{itemize}
  \item ヒルベルト空間を双対時空の位相空間で構成することで、量子状態を自然に表現できました。
  \item 位置演算子を双対時空の基底と並進セクターで定義し、スピノル（粒子のスピン）とつなげることで、量子力学と相対論が統一的に結びつきます。
  \item モーター \(M = R\, T\) が、回転・ブーストと並進を一つの演算子で扱い、局所化を美しく完成させます。
\end{itemize}

この理論の素晴らしいところは、「時空は連続的な一つのもの」という古い考えを捨て、粒子ごとに二つの時空と並進を持つという大胆なアイデアで、多くの難問を一度に解決する点です。重力や慣性も、ねじれ不整合（二つの時空のずれ）として説明できるため、量子重力の道が自然に開かれています。

初学者の方には少し難しい部分もあったかもしれませんが、二重時空理論（PGA版）は「二つの時空と並進が協力して、失われた情報を取り戻し、実際の移動を光的に実行する」というシンプルなイメージが核心です。この考え方が、みなさんの物理への興味をさらに深めてくれれば嬉しいです。

ご質問やもっと知りたいところがあれば、いつでも聞いてくださいね！ これからも一緒に学んでいきましょう。ありがとうございました！

\end{document}
